% 用 XeLaTeX 编译。默认生成题目版，定义 \AnswerEdition 后生成答案版。
% xelatex -jobname=三年级计算_题目版 三年级计算.tex
% xelatex -jobname=三年级计算_答案版 '\def\AnswerEdition{}% 用 XeLaTeX 编译。默认生成题目版，定义 \AnswerEdition 后生成答案版。
% xelatex -jobname=三年级计算_题目版 三年级计算.tex
% xelatex -jobname=三年级计算_答案版 '\def\AnswerEdition{}% 用 XeLaTeX 编译。默认生成题目版，定义 \AnswerEdition 后生成答案版。
% xelatex -jobname=三年级计算_题目版 三年级计算.tex
% xelatex -jobname=三年级计算_答案版 '\def\AnswerEdition{}% 用 XeLaTeX 编译。默认生成题目版，定义 \AnswerEdition 后生成答案版。
% xelatex -jobname=三年级计算_题目版 三年级计算.tex
% xelatex -jobname=三年级计算_答案版 '\def\AnswerEdition{}\input{三年级计算.tex}'
\documentclass[a4paper]{article}
\usepackage[left=17mm,right=17mm,top=15mm,bottom=13mm]{geometry}
\usepackage{fontspec}
\usepackage{unicode-math}
\usepackage{tikz}
\setmathfont{latinmodern-math.otf}
\IfFileExists{/usr/share/fonts/opentype/noto/NotoSerifCJK-Regular.ttc}{
  \newfontfamily\ChineseFont{NotoSerifCJK-Regular.ttc}[
    Path=/usr/share/fonts/opentype/noto/,
    UprightFeatures={FontIndex=2},
    BoldFont=NotoSerifCJK-Bold.ttc,
    BoldFeatures={FontIndex=2}
  ]
}{
  \IfFontExistsTF{NotoSerifCJKsc-Regular}{
    \newfontfamily\ChineseFont{NotoSerifCJKsc-Regular}[BoldFont=NotoSerifCJKsc-Bold]
  }{
    \newfontfamily\ChineseFont{FandolSong-Regular.otf}[BoldFont=FandolSong-Bold.otf]
  }
}
\newcommand{\Chinese}[1]{{\ChineseFont #1}}
\newcommand{\AnswerOnly}[1]{\ifdefined\AnswerEdition #1\fi}
\pagestyle{empty}
\setlength{\parindent}{0pt}

% 竖式的各个数位使用等宽坐标。第一个坐标对应个位。
\ExplSyntaxOn
\seq_new:N \l_worksheet_digit_seq
\NewDocumentCommand{\Digits}{mmm}{
  \seq_set_split:Nnn \l_worksheet_digit_seq {} {#3}
  \seq_reverse:N \l_worksheet_digit_seq
  \seq_map_indexed_inline:Nn \l_worksheet_digit_seq {
    \node[work~digit] at (\fp_eval:n{#1-(##1-1)*4.6},#2) {$##2$};
  }
}
\ExplSyntaxOff

\newcommand{\SectionTitle}[2]{
  \node[anchor=west,font=\fontsize{15}{18}\selectfont\bfseries]
    at (0,#1) {\Chinese{#2}};
}
\newcommand{\VerticalPrompt}[5]{
  \node[anchor=west,font=\fontsize{17}{20}\selectfont] at (#1,#2)
    {#3\(#4=\AnswerOnly{#5}\)\ifdefined\AnswerEdition\else\hspace{2mm}\rule{29mm}{0.35pt}\fi};
}
\newcommand{\Remainder}[2]{#1\mathord{\text{\Chinese{……}}}#2}

\newcommand{\Multiplication}[7]{
  \begin{scope}[shift={(#1,#2)}]
    \Digits{0}{0}{#3}
    \Digits{0}{6}{#4}
    \node[work digit] at (-24,6) {$\times$};
    \draw[arithmetic rule] (-26,9) -- (3,9);
    \Digits{0}{12.5}{#5}
    \Digits{0}{18.5}{#6}
    \draw[arithmetic rule] (-26,21.5) -- (3,21.5);
    \Digits{0}{25}{#7}
  \end{scope}
}

% 除法：先画商与除号，再逐行填入每步试除的积和带下的数字。
\newcommand{\DivisionTop}[2]{
  \Digits{0}{0}{#2}
  \Digits{0}{6}{#1}
  \node[work digit] at (-19,6) {$5$};
  \draw[arithmetic rule] (3.5,2.9) -- (-12.3,2.9)
    .. controls (-12.6,6) and (-13.6,8) .. (-15.1,9.2);
}
\newcommand{\DivisionLine}[3]{
  \draw[arithmetic rule] (#1,#3) -- (#2,#3);
}
\newcommand{\DivideFiveHundredEight}[2]{
  \begin{scope}[shift={(#1,#2)}]
    \DivisionTop{508}{101}
    \Digits{-9.2}{11.5}{5}
    \DivisionLine{-12}{-6.4}{14}
    \Digits{-4.6}{17}{0}
    \Digits{-4.6}{22.5}{0}
    \DivisionLine{-7.4}{-1.8}{25}
    \Digits{0}{28}{8}
    \Digits{0}{33.5}{5}
    \DivisionLine{-2.8}{2.8}{36}
    \Digits{0}{39}{3}
  \end{scope}
}
\newcommand{\DivideSixHundredThree}[2]{
  \begin{scope}[shift={(#1,#2)}]
    \DivisionTop{603}{120}
    \Digits{-9.2}{11.5}{5}
    \DivisionLine{-12}{-6.4}{14}
    \Digits{-4.6}{17}{10}
    \Digits{-4.6}{22.5}{10}
    \DivisionLine{-12}{-1.8}{25}
    \Digits{0}{28}{3}
    \Digits{0}{33.5}{0}
    \DivisionLine{-2.8}{2.8}{36}
    \Digits{0}{39}{3}
  \end{scope}
}

\newcommand{\CheckColumn}[6]{
  \begin{scope}[shift={(#1,#2)}]
    \Digits{0}{0}{#3}
    \Digits{0}{6}{#5}
    \node[work digit] at (-15,6) {$#4$};
    \draw[arithmetic rule] (-17,9) -- (3,9);
    \Digits{0}{12.5}{#6}
  \end{scope}
}
\newcommand{\DivisionCheck}[5]{
  \begin{scope}[shift={(#1,#2)}]
    \node[anchor=west,font=\fontsize{11}{14}\selectfont] at (-17,-7)
      {\Chinese{验算}};
    \CheckColumn{0}{0}{#3}{\times}{5}{#4}
    \CheckColumn{25}{0}{#4}{+}{3}{#5}
    \node[anchor=west,font=\fontsize{11}{14}\selectfont] at (-17,20)
      {\Chinese{余数}\enspace\(3<5\)};
  \end{scope}
}
\newcommand{\StepPrompt}[3]{
  \node[anchor=west,font=\fontsize{16}{20}\selectfont] at (#1+5,#2) {\(#3\)};
}
\newcommand{\Step}[3]{
  \node[anchor=west,font=\fontsize{16}{20}\selectfont] at (#1,#2) {\(=\,#3\)};
}

\begin{document}
\noindent
\begin{tikzpicture}[
    x=1mm,y=-1mm,
    every node/.style={inner sep=0pt,outer sep=0pt},
    work digit/.style={font=\fontsize{13}{15}\selectfont},
    arithmetic rule/.style={line width=0.5pt,line cap=round}
  ]
  \path[use as bounding box] (0,0) rectangle (176,258);
  \node[font=\fontsize{26}{31}\selectfont\bfseries] at (88,9)
    {\Chinese{三年级计算}};
  \ifdefined\AnswerEdition
    \node[font=\fontsize{12}{15}\selectfont] at (88,23) {\Chinese{答案版}};
  \else
    \node[anchor=west,font=\fontsize{12}{15}\selectfont] at (0,24)
      {\Chinese{班级}\hspace{2mm}\rule{29mm}{0.35pt}\hspace{14mm}%
       \Chinese{姓名}\hspace{2mm}\rule{29mm}{0.35pt}};
  \fi

  \SectionTitle{38}{竖式计算，带※的要验算}
  \VerticalPrompt{2}{51}{}{1008\times58}{58464}
  \VerticalPrompt{92}{51}{\Chinese{※}\hspace{1mm}}{508\div5}{\Remainder{101}{3}}
  % 照片左下的印刷题已划去，两版均采用手写订正题。
  \VerticalPrompt{2}{111}{\Chinese{※}\hspace{1mm}}{603\div5}{\Remainder{120}{3}}
  \VerticalPrompt{92}{111}{}{850\times32}{27200}
  \ifdefined\AnswerEdition
    \Multiplication{29}{64}{1008}{58}{8064}{50400}{58464}
    \DivideFiveHundredEight{113}{62}
    \DivisionCheck{141}{73}{101}{505}{508}
    \DivideSixHundredThree{23}{122}
    \DivisionCheck{51}{133}{120}{600}{603}
    \Multiplication{119}{124}{850}{32}{1700}{25500}{27200}
  \fi

  \SectionTitle{174}{脱式计算}
  \StepPrompt{0}{188}{540\div9\div3}
  \StepPrompt{59}{188}{1024\times2\div8}
  \StepPrompt{118}{188}{504\times(17+52)}
  \StepPrompt{0}{225}{(1105-455)\times12}
  \StepPrompt{59}{225}{99\times94+94}
  \StepPrompt{118}{225}{1050\div5+100}
  \ifdefined\AnswerEdition
    \Step{0}{196}{60\div3}
    \Step{0}{203}{20}
    \Step{59}{196}{2048\div8}
    \Step{59}{203}{256}
    \Step{118}{196}{504\times69}
    \Step{118}{203}{34776}
    \Step{0}{233}{650\times12}
    \Step{0}{240}{7800}
    \Step{59}{233}{(99+1)\times94}
    \Step{59}{240}{100\times94}
    \Step{59}{247}{9400}
    \Step{118}{233}{210+100}
    \Step{118}{240}{310}
  \fi
\end{tikzpicture}%
\par
\end{document}
'
\documentclass[a4paper]{article}
\usepackage[left=17mm,right=17mm,top=15mm,bottom=13mm]{geometry}
\usepackage{fontspec}
\usepackage{unicode-math}
\usepackage{tikz}
\setmathfont{latinmodern-math.otf}
\IfFileExists{/usr/share/fonts/opentype/noto/NotoSerifCJK-Regular.ttc}{
  \newfontfamily\ChineseFont{NotoSerifCJK-Regular.ttc}[
    Path=/usr/share/fonts/opentype/noto/,
    UprightFeatures={FontIndex=2},
    BoldFont=NotoSerifCJK-Bold.ttc,
    BoldFeatures={FontIndex=2}
  ]
}{
  \IfFontExistsTF{NotoSerifCJKsc-Regular}{
    \newfontfamily\ChineseFont{NotoSerifCJKsc-Regular}[BoldFont=NotoSerifCJKsc-Bold]
  }{
    \newfontfamily\ChineseFont{FandolSong-Regular.otf}[BoldFont=FandolSong-Bold.otf]
  }
}
\newcommand{\Chinese}[1]{{\ChineseFont #1}}
\newcommand{\AnswerOnly}[1]{\ifdefined\AnswerEdition #1\fi}
\pagestyle{empty}
\setlength{\parindent}{0pt}

% 竖式的各个数位使用等宽坐标。第一个坐标对应个位。
\ExplSyntaxOn
\seq_new:N \l_worksheet_digit_seq
\NewDocumentCommand{\Digits}{mmm}{
  \seq_set_split:Nnn \l_worksheet_digit_seq {} {#3}
  \seq_reverse:N \l_worksheet_digit_seq
  \seq_map_indexed_inline:Nn \l_worksheet_digit_seq {
    \node[work~digit] at (\fp_eval:n{#1-(##1-1)*4.6},#2) {$##2$};
  }
}
\ExplSyntaxOff

\newcommand{\SectionTitle}[2]{
  \node[anchor=west,font=\fontsize{15}{18}\selectfont\bfseries]
    at (0,#1) {\Chinese{#2}};
}
\newcommand{\VerticalPrompt}[5]{
  \node[anchor=west,font=\fontsize{17}{20}\selectfont] at (#1,#2)
    {#3\(#4=\AnswerOnly{#5}\)\ifdefined\AnswerEdition\else\hspace{2mm}\rule{29mm}{0.35pt}\fi};
}
\newcommand{\Remainder}[2]{#1\mathord{\text{\Chinese{……}}}#2}

\newcommand{\Multiplication}[7]{
  \begin{scope}[shift={(#1,#2)}]
    \Digits{0}{0}{#3}
    \Digits{0}{6}{#4}
    \node[work digit] at (-24,6) {$\times$};
    \draw[arithmetic rule] (-26,9) -- (3,9);
    \Digits{0}{12.5}{#5}
    \Digits{0}{18.5}{#6}
    \draw[arithmetic rule] (-26,21.5) -- (3,21.5);
    \Digits{0}{25}{#7}
  \end{scope}
}

% 除法：先画商与除号，再逐行填入每步试除的积和带下的数字。
\newcommand{\DivisionTop}[2]{
  \Digits{0}{0}{#2}
  \Digits{0}{6}{#1}
  \node[work digit] at (-19,6) {$5$};
  \draw[arithmetic rule] (3.5,2.9) -- (-12.3,2.9)
    .. controls (-12.6,6) and (-13.6,8) .. (-15.1,9.2);
}
\newcommand{\DivisionLine}[3]{
  \draw[arithmetic rule] (#1,#3) -- (#2,#3);
}
\newcommand{\DivideFiveHundredEight}[2]{
  \begin{scope}[shift={(#1,#2)}]
    \DivisionTop{508}{101}
    \Digits{-9.2}{11.5}{5}
    \DivisionLine{-12}{-6.4}{14}
    \Digits{-4.6}{17}{0}
    \Digits{-4.6}{22.5}{0}
    \DivisionLine{-7.4}{-1.8}{25}
    \Digits{0}{28}{8}
    \Digits{0}{33.5}{5}
    \DivisionLine{-2.8}{2.8}{36}
    \Digits{0}{39}{3}
  \end{scope}
}
\newcommand{\DivideSixHundredThree}[2]{
  \begin{scope}[shift={(#1,#2)}]
    \DivisionTop{603}{120}
    \Digits{-9.2}{11.5}{5}
    \DivisionLine{-12}{-6.4}{14}
    \Digits{-4.6}{17}{10}
    \Digits{-4.6}{22.5}{10}
    \DivisionLine{-12}{-1.8}{25}
    \Digits{0}{28}{3}
    \Digits{0}{33.5}{0}
    \DivisionLine{-2.8}{2.8}{36}
    \Digits{0}{39}{3}
  \end{scope}
}

\newcommand{\CheckColumn}[6]{
  \begin{scope}[shift={(#1,#2)}]
    \Digits{0}{0}{#3}
    \Digits{0}{6}{#5}
    \node[work digit] at (-15,6) {$#4$};
    \draw[arithmetic rule] (-17,9) -- (3,9);
    \Digits{0}{12.5}{#6}
  \end{scope}
}
\newcommand{\DivisionCheck}[5]{
  \begin{scope}[shift={(#1,#2)}]
    \node[anchor=west,font=\fontsize{11}{14}\selectfont] at (-17,-7)
      {\Chinese{验算}};
    \CheckColumn{0}{0}{#3}{\times}{5}{#4}
    \CheckColumn{25}{0}{#4}{+}{3}{#5}
    \node[anchor=west,font=\fontsize{11}{14}\selectfont] at (-17,20)
      {\Chinese{余数}\enspace\(3<5\)};
  \end{scope}
}
\newcommand{\StepPrompt}[3]{
  \node[anchor=west,font=\fontsize{16}{20}\selectfont] at (#1+5,#2) {\(#3\)};
}
\newcommand{\Step}[3]{
  \node[anchor=west,font=\fontsize{16}{20}\selectfont] at (#1,#2) {\(=\,#3\)};
}

\begin{document}
\noindent
\begin{tikzpicture}[
    x=1mm,y=-1mm,
    every node/.style={inner sep=0pt,outer sep=0pt},
    work digit/.style={font=\fontsize{13}{15}\selectfont},
    arithmetic rule/.style={line width=0.5pt,line cap=round}
  ]
  \path[use as bounding box] (0,0) rectangle (176,258);
  \node[font=\fontsize{26}{31}\selectfont\bfseries] at (88,9)
    {\Chinese{三年级计算}};
  \ifdefined\AnswerEdition
    \node[font=\fontsize{12}{15}\selectfont] at (88,23) {\Chinese{答案版}};
  \else
    \node[anchor=west,font=\fontsize{12}{15}\selectfont] at (0,24)
      {\Chinese{班级}\hspace{2mm}\rule{29mm}{0.35pt}\hspace{14mm}%
       \Chinese{姓名}\hspace{2mm}\rule{29mm}{0.35pt}};
  \fi

  \SectionTitle{38}{竖式计算，带※的要验算}
  \VerticalPrompt{2}{51}{}{1008\times58}{58464}
  \VerticalPrompt{92}{51}{\Chinese{※}\hspace{1mm}}{508\div5}{\Remainder{101}{3}}
  % 照片左下的印刷题已划去，两版均采用手写订正题。
  \VerticalPrompt{2}{111}{\Chinese{※}\hspace{1mm}}{603\div5}{\Remainder{120}{3}}
  \VerticalPrompt{92}{111}{}{850\times32}{27200}
  \ifdefined\AnswerEdition
    \Multiplication{29}{64}{1008}{58}{8064}{50400}{58464}
    \DivideFiveHundredEight{113}{62}
    \DivisionCheck{141}{73}{101}{505}{508}
    \DivideSixHundredThree{23}{122}
    \DivisionCheck{51}{133}{120}{600}{603}
    \Multiplication{119}{124}{850}{32}{1700}{25500}{27200}
  \fi

  \SectionTitle{174}{脱式计算}
  \StepPrompt{0}{188}{540\div9\div3}
  \StepPrompt{59}{188}{1024\times2\div8}
  \StepPrompt{118}{188}{504\times(17+52)}
  \StepPrompt{0}{225}{(1105-455)\times12}
  \StepPrompt{59}{225}{99\times94+94}
  \StepPrompt{118}{225}{1050\div5+100}
  \ifdefined\AnswerEdition
    \Step{0}{196}{60\div3}
    \Step{0}{203}{20}
    \Step{59}{196}{2048\div8}
    \Step{59}{203}{256}
    \Step{118}{196}{504\times69}
    \Step{118}{203}{34776}
    \Step{0}{233}{650\times12}
    \Step{0}{240}{7800}
    \Step{59}{233}{(99+1)\times94}
    \Step{59}{240}{100\times94}
    \Step{59}{247}{9400}
    \Step{118}{233}{210+100}
    \Step{118}{240}{310}
  \fi
\end{tikzpicture}%
\par
\end{document}
'
\documentclass[a4paper]{article}
\usepackage[left=17mm,right=17mm,top=15mm,bottom=13mm]{geometry}
\usepackage{fontspec}
\usepackage{unicode-math}
\usepackage{tikz}
\setmathfont{latinmodern-math.otf}
\IfFileExists{/usr/share/fonts/opentype/noto/NotoSerifCJK-Regular.ttc}{
  \newfontfamily\ChineseFont{NotoSerifCJK-Regular.ttc}[
    Path=/usr/share/fonts/opentype/noto/,
    UprightFeatures={FontIndex=2},
    BoldFont=NotoSerifCJK-Bold.ttc,
    BoldFeatures={FontIndex=2}
  ]
}{
  \IfFontExistsTF{NotoSerifCJKsc-Regular}{
    \newfontfamily\ChineseFont{NotoSerifCJKsc-Regular}[BoldFont=NotoSerifCJKsc-Bold]
  }{
    \newfontfamily\ChineseFont{FandolSong-Regular.otf}[BoldFont=FandolSong-Bold.otf]
  }
}
\newcommand{\Chinese}[1]{{\ChineseFont #1}}
\newcommand{\AnswerOnly}[1]{\ifdefined\AnswerEdition #1\fi}
\pagestyle{empty}
\setlength{\parindent}{0pt}

% 竖式的各个数位使用等宽坐标。第一个坐标对应个位。
\ExplSyntaxOn
\seq_new:N \l_worksheet_digit_seq
\NewDocumentCommand{\Digits}{mmm}{
  \seq_set_split:Nnn \l_worksheet_digit_seq {} {#3}
  \seq_reverse:N \l_worksheet_digit_seq
  \seq_map_indexed_inline:Nn \l_worksheet_digit_seq {
    \node[work~digit] at (\fp_eval:n{#1-(##1-1)*4.6},#2) {$##2$};
  }
}
\ExplSyntaxOff

\newcommand{\SectionTitle}[2]{
  \node[anchor=west,font=\fontsize{15}{18}\selectfont\bfseries]
    at (0,#1) {\Chinese{#2}};
}
\newcommand{\VerticalPrompt}[5]{
  \node[anchor=west,font=\fontsize{17}{20}\selectfont] at (#1,#2)
    {#3\(#4=\AnswerOnly{#5}\)\ifdefined\AnswerEdition\else\hspace{2mm}\rule{29mm}{0.35pt}\fi};
}
\newcommand{\Remainder}[2]{#1\mathord{\text{\Chinese{……}}}#2}

\newcommand{\Multiplication}[7]{
  \begin{scope}[shift={(#1,#2)}]
    \Digits{0}{0}{#3}
    \Digits{0}{6}{#4}
    \node[work digit] at (-24,6) {$\times$};
    \draw[arithmetic rule] (-26,9) -- (3,9);
    \Digits{0}{12.5}{#5}
    \Digits{0}{18.5}{#6}
    \draw[arithmetic rule] (-26,21.5) -- (3,21.5);
    \Digits{0}{25}{#7}
  \end{scope}
}

% 除法：先画商与除号，再逐行填入每步试除的积和带下的数字。
\newcommand{\DivisionTop}[2]{
  \Digits{0}{0}{#2}
  \Digits{0}{6}{#1}
  \node[work digit] at (-19,6) {$5$};
  \draw[arithmetic rule] (3.5,2.9) -- (-12.3,2.9)
    .. controls (-12.6,6) and (-13.6,8) .. (-15.1,9.2);
}
\newcommand{\DivisionLine}[3]{
  \draw[arithmetic rule] (#1,#3) -- (#2,#3);
}
\newcommand{\DivideFiveHundredEight}[2]{
  \begin{scope}[shift={(#1,#2)}]
    \DivisionTop{508}{101}
    \Digits{-9.2}{11.5}{5}
    \DivisionLine{-12}{-6.4}{14}
    \Digits{-4.6}{17}{0}
    \Digits{-4.6}{22.5}{0}
    \DivisionLine{-7.4}{-1.8}{25}
    \Digits{0}{28}{8}
    \Digits{0}{33.5}{5}
    \DivisionLine{-2.8}{2.8}{36}
    \Digits{0}{39}{3}
  \end{scope}
}
\newcommand{\DivideSixHundredThree}[2]{
  \begin{scope}[shift={(#1,#2)}]
    \DivisionTop{603}{120}
    \Digits{-9.2}{11.5}{5}
    \DivisionLine{-12}{-6.4}{14}
    \Digits{-4.6}{17}{10}
    \Digits{-4.6}{22.5}{10}
    \DivisionLine{-12}{-1.8}{25}
    \Digits{0}{28}{3}
    \Digits{0}{33.5}{0}
    \DivisionLine{-2.8}{2.8}{36}
    \Digits{0}{39}{3}
  \end{scope}
}

\newcommand{\CheckColumn}[6]{
  \begin{scope}[shift={(#1,#2)}]
    \Digits{0}{0}{#3}
    \Digits{0}{6}{#5}
    \node[work digit] at (-15,6) {$#4$};
    \draw[arithmetic rule] (-17,9) -- (3,9);
    \Digits{0}{12.5}{#6}
  \end{scope}
}
\newcommand{\DivisionCheck}[5]{
  \begin{scope}[shift={(#1,#2)}]
    \node[anchor=west,font=\fontsize{11}{14}\selectfont] at (-17,-7)
      {\Chinese{验算}};
    \CheckColumn{0}{0}{#3}{\times}{5}{#4}
    \CheckColumn{25}{0}{#4}{+}{3}{#5}
    \node[anchor=west,font=\fontsize{11}{14}\selectfont] at (-17,20)
      {\Chinese{余数}\enspace\(3<5\)};
  \end{scope}
}
\newcommand{\StepPrompt}[3]{
  \node[anchor=west,font=\fontsize{16}{20}\selectfont] at (#1+5,#2) {\(#3\)};
}
\newcommand{\Step}[3]{
  \node[anchor=west,font=\fontsize{16}{20}\selectfont] at (#1,#2) {\(=\,#3\)};
}

\begin{document}
\noindent
\begin{tikzpicture}[
    x=1mm,y=-1mm,
    every node/.style={inner sep=0pt,outer sep=0pt},
    work digit/.style={font=\fontsize{13}{15}\selectfont},
    arithmetic rule/.style={line width=0.5pt,line cap=round}
  ]
  \path[use as bounding box] (0,0) rectangle (176,258);
  \node[font=\fontsize{26}{31}\selectfont\bfseries] at (88,9)
    {\Chinese{三年级计算}};
  \ifdefined\AnswerEdition
    \node[font=\fontsize{12}{15}\selectfont] at (88,23) {\Chinese{答案版}};
  \else
    \node[anchor=west,font=\fontsize{12}{15}\selectfont] at (0,24)
      {\Chinese{班级}\hspace{2mm}\rule{29mm}{0.35pt}\hspace{14mm}%
       \Chinese{姓名}\hspace{2mm}\rule{29mm}{0.35pt}};
  \fi

  \SectionTitle{38}{竖式计算，带※的要验算}
  \VerticalPrompt{2}{51}{}{1008\times58}{58464}
  \VerticalPrompt{92}{51}{\Chinese{※}\hspace{1mm}}{508\div5}{\Remainder{101}{3}}
  % 照片左下的印刷题已划去，两版均采用手写订正题。
  \VerticalPrompt{2}{111}{\Chinese{※}\hspace{1mm}}{603\div5}{\Remainder{120}{3}}
  \VerticalPrompt{92}{111}{}{850\times32}{27200}
  \ifdefined\AnswerEdition
    \Multiplication{29}{64}{1008}{58}{8064}{50400}{58464}
    \DivideFiveHundredEight{113}{62}
    \DivisionCheck{141}{73}{101}{505}{508}
    \DivideSixHundredThree{23}{122}
    \DivisionCheck{51}{133}{120}{600}{603}
    \Multiplication{119}{124}{850}{32}{1700}{25500}{27200}
  \fi

  \SectionTitle{174}{脱式计算}
  \StepPrompt{0}{188}{540\div9\div3}
  \StepPrompt{59}{188}{1024\times2\div8}
  \StepPrompt{118}{188}{504\times(17+52)}
  \StepPrompt{0}{225}{(1105-455)\times12}
  \StepPrompt{59}{225}{99\times94+94}
  \StepPrompt{118}{225}{1050\div5+100}
  \ifdefined\AnswerEdition
    \Step{0}{196}{60\div3}
    \Step{0}{203}{20}
    \Step{59}{196}{2048\div8}
    \Step{59}{203}{256}
    \Step{118}{196}{504\times69}
    \Step{118}{203}{34776}
    \Step{0}{233}{650\times12}
    \Step{0}{240}{7800}
    \Step{59}{233}{(99+1)\times94}
    \Step{59}{240}{100\times94}
    \Step{59}{247}{9400}
    \Step{118}{233}{210+100}
    \Step{118}{240}{310}
  \fi
\end{tikzpicture}%
\par
\end{document}
'
\documentclass[a4paper]{article}
\usepackage[left=17mm,right=17mm,top=15mm,bottom=13mm]{geometry}
\usepackage{fontspec}
\usepackage{unicode-math}
\usepackage{tikz}
\setmathfont{latinmodern-math.otf}
\IfFileExists{/usr/share/fonts/opentype/noto/NotoSerifCJK-Regular.ttc}{
  \newfontfamily\ChineseFont{NotoSerifCJK-Regular.ttc}[
    Path=/usr/share/fonts/opentype/noto/,
    UprightFeatures={FontIndex=2},
    BoldFont=NotoSerifCJK-Bold.ttc,
    BoldFeatures={FontIndex=2}
  ]
}{
  \IfFontExistsTF{NotoSerifCJKsc-Regular}{
    \newfontfamily\ChineseFont{NotoSerifCJKsc-Regular}[BoldFont=NotoSerifCJKsc-Bold]
  }{
    \newfontfamily\ChineseFont{FandolSong-Regular.otf}[BoldFont=FandolSong-Bold.otf]
  }
}
\newcommand{\Chinese}[1]{{\ChineseFont #1}}
\newcommand{\AnswerOnly}[1]{\ifdefined\AnswerEdition #1\fi}
\pagestyle{empty}
\setlength{\parindent}{0pt}

% 竖式的各个数位使用等宽坐标。第一个坐标对应个位。
\ExplSyntaxOn
\seq_new:N \l_worksheet_digit_seq
\NewDocumentCommand{\Digits}{mmm}{
  \seq_set_split:Nnn \l_worksheet_digit_seq {} {#3}
  \seq_reverse:N \l_worksheet_digit_seq
  \seq_map_indexed_inline:Nn \l_worksheet_digit_seq {
    \node[work~digit] at (\fp_eval:n{#1-(##1-1)*4.6},#2) {$##2$};
  }
}
\ExplSyntaxOff

\newcommand{\SectionTitle}[2]{
  \node[anchor=west,font=\fontsize{15}{18}\selectfont\bfseries]
    at (0,#1) {\Chinese{#2}};
}
\newcommand{\VerticalPrompt}[5]{
  \node[anchor=west,font=\fontsize{17}{20}\selectfont] at (#1,#2)
    {#3\(#4=\AnswerOnly{#5}\)\ifdefined\AnswerEdition\else\hspace{2mm}\rule{29mm}{0.35pt}\fi};
}
\newcommand{\Remainder}[2]{#1\mathord{\text{\Chinese{……}}}#2}

\newcommand{\Multiplication}[7]{
  \begin{scope}[shift={(#1,#2)}]
    \Digits{0}{0}{#3}
    \Digits{0}{6}{#4}
    \node[work digit] at (-24,6) {$\times$};
    \draw[arithmetic rule] (-26,9) -- (3,9);
    \Digits{0}{12.5}{#5}
    \Digits{0}{18.5}{#6}
    \draw[arithmetic rule] (-26,21.5) -- (3,21.5);
    \Digits{0}{25}{#7}
  \end{scope}
}

% 除法：先画商与除号，再逐行填入每步试除的积和带下的数字。
\newcommand{\DivisionTop}[2]{
  \Digits{0}{0}{#2}
  \Digits{0}{6}{#1}
  \node[work digit] at (-19,6) {$5$};
  \draw[arithmetic rule] (3.5,2.9) -- (-12.3,2.9)
    .. controls (-12.6,6) and (-13.6,8) .. (-15.1,9.2);
}
\newcommand{\DivisionLine}[3]{
  \draw[arithmetic rule] (#1,#3) -- (#2,#3);
}
\newcommand{\DivideFiveHundredEight}[2]{
  \begin{scope}[shift={(#1,#2)}]
    \DivisionTop{508}{101}
    \Digits{-9.2}{11.5}{5}
    \DivisionLine{-12}{-6.4}{14}
    \Digits{-4.6}{17}{0}
    \Digits{-4.6}{22.5}{0}
    \DivisionLine{-7.4}{-1.8}{25}
    \Digits{0}{28}{8}
    \Digits{0}{33.5}{5}
    \DivisionLine{-2.8}{2.8}{36}
    \Digits{0}{39}{3}
  \end{scope}
}
\newcommand{\DivideSixHundredThree}[2]{
  \begin{scope}[shift={(#1,#2)}]
    \DivisionTop{603}{120}
    \Digits{-9.2}{11.5}{5}
    \DivisionLine{-12}{-6.4}{14}
    \Digits{-4.6}{17}{10}
    \Digits{-4.6}{22.5}{10}
    \DivisionLine{-12}{-1.8}{25}
    \Digits{0}{28}{3}
    \Digits{0}{33.5}{0}
    \DivisionLine{-2.8}{2.8}{36}
    \Digits{0}{39}{3}
  \end{scope}
}

\newcommand{\CheckColumn}[6]{
  \begin{scope}[shift={(#1,#2)}]
    \Digits{0}{0}{#3}
    \Digits{0}{6}{#5}
    \node[work digit] at (-15,6) {$#4$};
    \draw[arithmetic rule] (-17,9) -- (3,9);
    \Digits{0}{12.5}{#6}
  \end{scope}
}
\newcommand{\DivisionCheck}[5]{
  \begin{scope}[shift={(#1,#2)}]
    \node[anchor=west,font=\fontsize{11}{14}\selectfont] at (-17,-7)
      {\Chinese{验算}};
    \CheckColumn{0}{0}{#3}{\times}{5}{#4}
    \CheckColumn{25}{0}{#4}{+}{3}{#5}
    \node[anchor=west,font=\fontsize{11}{14}\selectfont] at (-17,20)
      {\Chinese{余数}\enspace\(3<5\)};
  \end{scope}
}
\newcommand{\StepPrompt}[3]{
  \node[anchor=west,font=\fontsize{16}{20}\selectfont] at (#1+5,#2) {\(#3\)};
}
\newcommand{\Step}[3]{
  \node[anchor=west,font=\fontsize{16}{20}\selectfont] at (#1,#2) {\(=\,#3\)};
}

\begin{document}
\noindent
\begin{tikzpicture}[
    x=1mm,y=-1mm,
    every node/.style={inner sep=0pt,outer sep=0pt},
    work digit/.style={font=\fontsize{13}{15}\selectfont},
    arithmetic rule/.style={line width=0.5pt,line cap=round}
  ]
  \path[use as bounding box] (0,0) rectangle (176,258);
  \node[font=\fontsize{26}{31}\selectfont\bfseries] at (88,9)
    {\Chinese{三年级计算}};
  \ifdefined\AnswerEdition
    \node[font=\fontsize{12}{15}\selectfont] at (88,23) {\Chinese{答案版}};
  \else
    \node[anchor=west,font=\fontsize{12}{15}\selectfont] at (0,24)
      {\Chinese{班级}\hspace{2mm}\rule{29mm}{0.35pt}\hspace{14mm}%
       \Chinese{姓名}\hspace{2mm}\rule{29mm}{0.35pt}};
  \fi

  \SectionTitle{38}{竖式计算，带※的要验算}
  \VerticalPrompt{2}{51}{}{1008\times58}{58464}
  \VerticalPrompt{92}{51}{\Chinese{※}\hspace{1mm}}{508\div5}{\Remainder{101}{3}}
  % 照片左下的印刷题已划去，两版均采用手写订正题。
  \VerticalPrompt{2}{111}{\Chinese{※}\hspace{1mm}}{603\div5}{\Remainder{120}{3}}
  \VerticalPrompt{92}{111}{}{850\times32}{27200}
  \ifdefined\AnswerEdition
    \Multiplication{29}{64}{1008}{58}{8064}{50400}{58464}
    \DivideFiveHundredEight{113}{62}
    \DivisionCheck{141}{73}{101}{505}{508}
    \DivideSixHundredThree{23}{122}
    \DivisionCheck{51}{133}{120}{600}{603}
    \Multiplication{119}{124}{850}{32}{1700}{25500}{27200}
  \fi

  \SectionTitle{174}{脱式计算}
  \StepPrompt{0}{188}{540\div9\div3}
  \StepPrompt{59}{188}{1024\times2\div8}
  \StepPrompt{118}{188}{504\times(17+52)}
  \StepPrompt{0}{225}{(1105-455)\times12}
  \StepPrompt{59}{225}{99\times94+94}
  \StepPrompt{118}{225}{1050\div5+100}
  \ifdefined\AnswerEdition
    \Step{0}{196}{60\div3}
    \Step{0}{203}{20}
    \Step{59}{196}{2048\div8}
    \Step{59}{203}{256}
    \Step{118}{196}{504\times69}
    \Step{118}{203}{34776}
    \Step{0}{233}{650\times12}
    \Step{0}{240}{7800}
    \Step{59}{233}{(99+1)\times94}
    \Step{59}{240}{100\times94}
    \Step{59}{247}{9400}
    \Step{118}{233}{210+100}
    \Step{118}{240}{310}
  \fi
\end{tikzpicture}%
\par
\end{document}
